\chapter{An Execution Is an Event Stream}
\label{ch:execution}
\label{ch:events}

\index{event stream}
\index{execution}
A protocol describes every conversation that may happen; an execution is one that
did. LeanFM represents an execution as the stream of its events, and keeps no
other state: roles keep none of their own that a requirement may depend on. What
a role knows at any moment is the part of the stream it has seen, and every
question about a running system is a question about the stream.

\section{Events}

\index{cause}
Each event is one message sent, or one thing a role did on its own. It carries
its session, its sender and receiver, its label, and the event that caused it.
Generated from \texttt{PlaceOrder} with its seed
6, one session:

\begin{lstlisting}[language={},basicstyle=\ttfamily\footnotesize,breaklines=true]
id 0  t=4 ms  CreateOrder: Client -> Shop
id 1  t=5 ms  Quote: Shop -> Client cause 0
id 2  t=6 ms  Pay: Client -> Shop cause 1
id 3  t=7 ms  Charge: Shop -> Bank cause 2
id 4  t=8 ms  Charged: Bank -> Shop cause 3
id 5  t=9 ms  Receipt: Shop -> Client cause 4
\end{lstlisting}


The cause links form the conversation's spine: each answer points at its
question, and the chain from \texttt{Receipt} back to \texttt{CreateOrder} is the
whole call. A real program logs the same shape, one line per message, and the
code generated from a protocol logs it automatically (Chapter~\ref{ch:generated}).

\section{Conformance}

\index{conformance}
A log conforms when every session in it is a sentence of the protocol: each event
allowed at its point, naming the cause the protocol names, and each session ending
where the protocol can end. Here the shop charges the card before the client has
paid:

\begin{lstlisting}[language={},basicstyle=\ttfamily\footnotesize,breaklines=true]
{"from":"Client","id":"e1","label":"CreateOrder","session":"42","to":"Shop"}
{"cause":"e1","from":"Shop","id":"e2","label":"Quote","session":"42","to":"Client"}
{"cause":"e2","from":"Shop","id":"e3","label":"Charge","session":"42","to":"Bank"}
\end{lstlisting}


\begin{lstlisting}[language={},basicstyle=\ttfamily\footnotesize,breaklines=true]
PlaceOrder: 1 of 1 sessions break the protocol; first: session 42, event e3: Charge from Shop to Bank is not allowed here; the protocol allows Pay from Client to Shop
session 42, event e3: Charge from Shop to Bank is not allowed here; the protocol allows Pay from Client to Shop
\end{lstlisting}


The verdict names the first event that breaks the protocol and what the protocol
allowed instead. Asked after any prefix, LeanFM also says what may happen next
and whether the session may end there; that is how a walk through one session is
driven, one click at a time.

\section{State lives in the stream}

\index{reward vector}
\index{view}
For questions about value and promises, events are keyed by a unique timestamp,
and two links point back in time: \texttt{prev}, the event before this one in its
thread, and \texttt{cause}, the event that made it happen. Both must name earlier
events of the same stream. A role's \emph{view} is its own events and everything
in their causal past: exactly what it could have learned.

Rewards follow the same rule. A reward vector, the balance of every role in every
unit, is written only on an event where it changes, and writing one equal to the
vector it inherits is refused. The vector at any event is found by walking
\texttt{prev} back to the latest one written; finding none means zero. A
concurrent event's vector is never taken, because it is not on the walk. In this
stream each ask moves the client's spend and the third ask is never answered:

\lstinputlisting[language={},basicstyle=\ttfamily\scriptsize,breaklines=true]{generated/sla-execution.json}

\section{Promises on a stream}

\index{observation period}
A \texttt{require} promise is judged on the stream as it was at design time, with
causality in place of probability. An entry hits when a target lies in its causal
future within the deadline; an answer that names no cause does not count, however
soon it comes. An entry whose deadline outlives the stream is pending, not missed.
Each window gets an exact one-sided binomial test, but only once it has closed;
until then it is open:

\lstinputlisting[language={},basicstyle=\ttfamily\scriptsize,breaklines=true]{generated/sla-execution.txt}

One miss in four would break a promise of 99 in 100 if it were judged at once.
The 300-entry period has not closed, so the promise is open. Over 300 closed
entries, 300 hits is met at 95/100 confidence ($0.99^{300} \approx 0.049$), 299
is undecided, and 290 is violated.

\section{Streams from the protocol}

\index{simulation}
The same probabilities that judge a promise can generate streams. A seeded
generator walks the protocol, drawing each choice, delay and amount by its exact
probability, so the same seed gives the same stream. Every generated session is a
walk of the protocol, so it conforms by construction; fed back as an execution,
it gives each session's realized rewards, and over many sessions their mean
approaches the expected value of Chapter~\ref{ch:assumptions}. Generated streams
are how an implementation is exercised before real traffic exists, and real
streams are how its assumptions are corrected once it does.
